\documentclass[11pt,a4paper]{article}

\usepackage[utf8]{inputenc}
\usepackage[T1]{fontenc}
\usepackage[ngerman]{babel}

\usepackage{amsmath,amssymb,amsthm,mathtools}
\usepackage[a4paper,margin=2.5cm]{geometry}
\usepackage{lmodern}
\usepackage{microtype}
\usepackage{hyperref}
\usepackage{enumitem}

\title{Ein \(E\!\)-/\(ABC\)-Invarianter aus der geschichteten Radikalzerlegung}
\author{Kollaborative Ausarbeitung}
\date{\today}

\newtheorem{satz}{Satz}
\newtheorem{lemma}[satz]{Lemma}
\newtheorem{korollar}[satz]{Korollar}
\theoremstyle{definition}
\newtheorem{definition}[satz]{Definition}
\newtheorem{beispiel}[satz]{Beispiel}
\theoremstyle{remark}
\newtheorem{bemerkung}[satz]{Bemerkung}

\newcommand{\N}{\mathbb{N}}
\newcommand{\Q}{\mathbb{Q}}
\newcommand{\ggT}{\operatorname{ggT}}
\newcommand{\Typ}{\operatorname{Typ}}
\newcommand{\Rest}{\mathcal{R}}
\newcommand{\rad}{\operatorname{rad}}

\begin{document}
	\maketitle
	
	\begin{abstract}
		Wir formulieren eine kanonische Zerlegung natürlicher Zahlen in einen \(2\)-Anteil, einen
		\(3\)-Anteil, einen neutralen \(E\)-Anteil und einen reduzierten \(ABC\)-Anteil.
		Aus dieser Zerlegung wird ein rationaler \(E\!\)-/\(ABC\)-Invarianter konstruiert.
		Dabei ist nicht die Zahl selbst ein Quotient ihres \(ABC\)- und \(E\)-Anteils,
		sondern der Quotient ist ein aus der Zahl eindeutig abgeleiteter Strukturinvariant.
	\end{abstract}
	
	\section{Ausgangslage}
	
	Wir schreiben für \(n\in\N\)
	\[
	n = 2^\alpha 3^\beta m,
	\qquad
	\alpha := v_2(n),\quad \beta := v_3(n),\quad \ggT(m,6)=1.
	\]
	Für den zu \(6\) teilerfremden Anteil \(m\) betrachten wir die geschichtete Radikalzerlegung
	\[
	m=\prod_{k=1}^h R_k,
	\qquad
	R_k:=\prod_{\substack{p\mid m\\ v_p(m)\ge k}} p.
	\]
	Jede Schicht \(R_k\) ist quadratfrei und zu \(6\) teilerfremd.
	
	Die Primzahlen \(p\ge 5\) zerfallen modulo \(12\) in die vier Familien
	\[
	E:=\{p:\ p\equiv 1 \pmod{12}\},\qquad
	A:=\{p:\ p\equiv 5 \pmod{12}\},
	\]
	\[
	B:=\{p:\ p\equiv 7 \pmod{12}\},\qquad
	C:=\{p:\ p\equiv 11 \pmod{12}\}.
	\]
	
	\section{Kanonische Repräsentanten der reduzierten Typen}
	
	\begin{definition}[Sieben reduzierte Typen]
		Wir setzen
		\[
		\mathcal{T}:=\{E,A,B,C,AB,AC,BC\}.
		\]
		Jede Radikalschicht \(R_k\) besitzt genau einen reduzierten Typ
		\[
		\Typ(R_k)\in\mathcal{T}.
		\]
	\end{definition}
	
	\begin{definition}[Kanonische numerische Repräsentanten]
		Wir definieren eine Abbildung
		\[
		\rho:\mathcal{T}\to\N
		\]
		durch
		\[
		\rho(E)=1,\qquad \rho(A)=5,\qquad \rho(B)=7,\qquad \rho(C)=11,
		\]
		\[
		\rho(AB)=5\cdot 7=35,\qquad
		\rho(AC)=5\cdot 11=55,\qquad
		\rho(BC)=7\cdot 11=77.
		\]
	\end{definition}
	
	\begin{bemerkung}
		Die Wahl von \(\rho\) ist eine Konvention. Sie macht die reduzierten Typen numerisch fassbar.
		Andere feste Repräsentantensysteme wären ebenfalls möglich. Die nachfolgenden Invarianten sind
		also relativ zu dieser kanonischen Wahl definiert.
	\end{bemerkung}
	
	\section{Kanonische \(E\)-/Rest-Zerlegung jeder Schicht}
	
	\begin{definition}[Schichtweiser \(E\)-Anteil]
		Für eine Radikalschicht \(R_k\) sei
		\[
		E_k:=\prod_{\substack{p\mid R_k\\ p\equiv 1\;(\mathrm{mod}\;12)}} p
		\]
		das Produkt aller Primteiler der Schicht aus der Familie \(E\).
	\end{definition}
	
	\begin{definition}[Schichtweiser reduzierter Restfaktor]
		Für eine Radikalschicht \(R_k\) definieren wir
		\[
		X_k:=\rho\bigl(\Typ(R_k)\bigr).
		\]
	\end{definition}
	
	\begin{bemerkung}
		Der Faktor \(X_k\) ist kein echter Teiler von \(R_k\), sondern ein numerischer Repräsentant
		des reduzierten Typs der Schicht. Er speichert also nicht die ganze arithmetische Information
		der Schicht, sondern nur ihre reduzierte \(ABC\)-Signatur.
	\end{bemerkung}
	
	\section{Globale \(E\)- und \(ABC\)-Invarianten}
	
	\begin{definition}[Globaler \(E\)-Anteil]
		Wir definieren den globalen \(E\)-Anteil von \(n\) durch
		\[
		E_{\mathrm{tot}}(n):=\prod_{k=1}^h E_k.
		\]
	\end{definition}
	
	\begin{definition}[Globaler reduzierter \(ABC\)-Anteil]
		Wir definieren den globalen reduzierten \(ABC\)-Anteil von \(n\) durch
		\[
		ABC_{\mathrm{tot}}(n):=\prod_{k=1}^h X_k
		=\prod_{k=1}^h \rho\bigl(\Typ(R_k)\bigr).
		\]
	\end{definition}
	
	\begin{definition}[\(E\!\)-/\(ABC\)-Invarianter]
		Der zu \(n\) gehörige rationale \(E\!\)-/\(ABC\)-Invarianter ist
		\[
		I_{E/ABC}(n):=\frac{ABC_{\mathrm{tot}}(n)}{E_{\mathrm{tot}}(n)}\in\Q_{>0}.
		\]
	\end{definition}
	
	\section{Eindeutigkeit}
	
	\begin{satz}[Kanonische Bestimmtheit]
		Zu jeder natürlichen Zahl \(n\in\N\) sind die Größen
		\[
		E_{\mathrm{tot}}(n),\qquad ABC_{\mathrm{tot}}(n),\qquad I_{E/ABC}(n)
		\]
		eindeutig bestimmt.
	\end{satz}
	
	\begin{proof}
		Die Zahlen \(v_2(n)\) und \(v_3(n)\) sind eindeutig bestimmt. Damit ist der zu \(6\) teilerfremde
		Anteil
		\[
		m=\frac{n}{2^{v_2(n)}3^{v_3(n)}}
		\]
		eindeutig festgelegt. Die Radikalschichten \(R_1,\dots,R_h\) sind dadurch ebenfalls eindeutig.
		
		Für jede Schicht \(R_k\) ist das Produkt \(E_k\) aller \(E\)-Primteiler eindeutig definiert.
		Zugleich ist der reduzierte Typ \(\Typ(R_k)\) eindeutig bestimmt, also auch sein numerischer
		Repräsentant \(X_k=\rho(\Typ(R_k))\). Folglich sind die Produkte
		\[
		E_{\mathrm{tot}}(n)=\prod_{k=1}^h E_k
		\qquad\text{und}\qquad
		ABC_{\mathrm{tot}}(n)=\prod_{k=1}^h X_k
		\]
		eindeutig. Damit ist auch der Quotient
		\[
		I_{E/ABC}(n)=\frac{ABC_{\mathrm{tot}}(n)}{E_{\mathrm{tot}}(n)}
		\]
		eindeutig bestimmt.
	\end{proof}
	
	\section{Die richtige Interpretation}
	
	\begin{satz}[Produktzerlegung mit rationalem Strukturinvarianten]
		Jede natürliche Zahl \(n\) besitzt eine kanonische Zerlegung der Form
		\[
		n = 2^{v_2(n)}\,3^{v_3(n)}\,m,
		\qquad
		m=\prod_{k=1}^h R_k,
		\]
		und daraus einen eindeutig bestimmten rationalen Strukturinvarianten
		\[
		I_{E/ABC}(n)=\frac{ABC_{\mathrm{tot}}(n)}{E_{\mathrm{tot}}(n)}.
		\]
		Im Allgemeinen gilt jedoch nicht
		\[
		n=\frac{ABC_{\mathrm{tot}}(n)}{E_{\mathrm{tot}}(n)}.
		\]
	\end{satz}
	
	\begin{proof}
		Die Existenz und Eindeutigkeit der Zerlegung sowie des Invarianten wurde bereits gezeigt.
		Dass im Allgemeinen
		\[
		n\neq \frac{ABC_{\mathrm{tot}}(n)}{E_{\mathrm{tot}}(n)}
		\]
		gilt, sieht man schon an einem einfachen Beispiel: Für \(n=5\cdot 7\cdot 13\) ist
		\[
		E_{\mathrm{tot}}(n)=13,\qquad ABC_{\mathrm{tot}}(n)=35,
		\]
		also
		\[
		I_{E/ABC}(n)=\frac{35}{13},
		\]
		während
		\[
		n=455.
		\]
	\end{proof}
	
	\begin{bemerkung}
		Die richtige Aussage lautet also nicht:
		\[
		n=\frac{\text{\(ABC\)-Anteil}}{\text{\(E\)-Anteil}},
		\]
		sondern:
		\[
		n \longmapsto \frac{\text{kanonischer reduzierter \(ABC\)-Anteil}}{\text{kanonischer \(E\)-Anteil}}
		\]
		definiert einen rationalen Invarianten.
	\end{bemerkung}
	
	\section{Alternative symbolische Schreibweise}
	
	\begin{definition}[Symbolische Signaturform]
		Mit
		\[
		\Rest(n)=\bigl(\Typ(R_1),\dots,\Typ(R_h)\bigr)
		\]
		kann man \(n\) symbolisch auch in der Form
		\[
		n \sim \bigl(v_2(n),\,v_3(n);\;E_{\mathrm{tot}}(n),\,\Rest(n)\bigr)
		\]
		auffassen.
		\]
	\end{definition}
	
	\begin{bemerkung}
		Diese symbolische Form ist oft konzeptionell klarer als der Quotient \(I_{E/ABC}(n)\), weil sie
		den gesamten reduzierten Rest als Wort über dem Alphabet
		\[
		\{E,A,B,C,AB,AC,BC\}
		\]
		sichtbar macht, statt ihn in einer einzelnen rationalen Zahl zusammenzufassen.
	\end{bemerkung}
	
	\section{Beispiele}
	
	\begin{beispiel}
		Für
		\[
		n=5\cdot 7\cdot 13
		\]
		gibt es nur eine Schicht
		\[
		R_1=5\cdot 7\cdot 13.
		\]
		Dann ist
		\[
		E_1=13,\qquad \Typ(R_1)=AB,\qquad X_1=\rho(AB)=35.
		\]
		Also
		\[
		E_{\mathrm{tot}}(n)=13,\qquad
		ABC_{\mathrm{tot}}(n)=35,
		\qquad
		I_{E/ABC}(n)=\frac{35}{13}.
		\]
	\end{beispiel}
	
	\begin{beispiel}
		Für
		\[
		n=2^2\cdot 3\cdot 5^3\cdot 7^2\cdot 13
		\]
		ist der zu \(6\) teilerfremde Anteil
		\[
		m=5^3\cdot 7^2\cdot 13.
		\]
		Die Radikalschichten sind
		\[
		R_1=5\cdot 7\cdot 13,\qquad
		R_2=5\cdot 7,\qquad
		R_3=5.
		\]
		Daraus folgt
		\[
		E_1=13,\qquad E_2=1,\qquad E_3=1,
		\]
		sowie
		\[
		\Typ(R_1)=AB,\qquad \Typ(R_2)=AB,\qquad \Typ(R_3)=A.
		\]
		Also
		\[
		ABC_{\mathrm{tot}}(n)=35\cdot 35\cdot 5=6125,
		\]
		\[
		E_{\mathrm{tot}}(n)=13,
		\]
		und damit
		\[
		I_{E/ABC}(n)=\frac{6125}{13}.
		\]
	\end{beispiel}
	
	\begin{beispiel}
		Für die Primzahl
		\[
		n=13
		\]
		gilt
		\[
		R_1=13,\qquad E_1=13,\qquad \Typ(R_1)=E,\qquad X_1=\rho(E)=1.
		\]
		Also
		\[
		I_{E/ABC}(13)=\frac{1}{13}.
		\]
	\end{beispiel}
	
	\section{Schluss}
	
	Jede natürliche Zahl besitzt nach Abspaltung des \(2\)- und \(3\)-Anteils eine eindeutige
	geschichtete Radikalzerlegung. Aus dieser ergibt sich kanonisch
	
	\begin{enumerate}[label=\arabic*.]
		\item ein globaler \(E\)-Anteil \(E_{\mathrm{tot}}(n)\),
		\item ein globaler reduzierter \(ABC\)-Anteil \(ABC_{\mathrm{tot}}(n)\),
		\item sowie der rationale Strukturinvariant
		\[
		I_{E/ABC}(n)=\frac{ABC_{\mathrm{tot}}(n)}{E_{\mathrm{tot}}(n)}.
		\]
	\end{enumerate}
	
	Damit erhält man keine Bruchdarstellung der Zahl selbst, wohl aber eine eindeutig definierte
	rationale Signatur der Zahl relativ zu ihrer \(E\)-/\(ABC\)-Struktur.
	
\end{document}